% Generated draft: requires mathematical and source review.
\documentclass[11pt]{article}
\usepackage[margin=25mm]{geometry}
\usepackage{fontspec}
\setmainfont{Latin Modern Roman}
\usepackage{amsmath,amssymb,mathtools,unicode-math}
\setmathfont{Latin Modern Math}
\usepackage{enumitem,needspace,xurl,hyperref}
\setlength{\parindent}{0pt}
\setlength{\parskip}{5pt}
\setlength{\emergencystretch}{3em}
\newcommand{\N}{\mathbb{N}}
\newcommand{\Z}{\mathbb{Z}}
\newcommand{\Q}{\mathbb{Q}}
\newcommand{\R}{\mathbb{R}}
\newcommand{\C}{\mathbb{C}}
\newcommand{\F}{\mathbb{F}}
\newcommand{\A}{\mathbb{A}}
\newcommand{\PP}{\mathbb{P}}
\newcommand{\E}{\mathbb{E}}
\newcommand{\eps}{\varepsilon}
\newcommand{\dd}{\,\mathrm{d}}
\begin{document}
\section*{UnsolvedMath ID: 30003521}
\textbf{Source problem code:} catalogue entry
\subsection{Definitions and mathematical statement}
The supplied problem concerns nonlinear wave systems on periodic graphs, with nonlinearities containing quadratic terms, and a nonlinear Schr\"odinger (NLS) approximation. The background refers to quantum graphs, suggesting a metric graph interpretation: functions are defined on edges parametrized by intervals, and differential equations hold along those edges. The particular graph, its periodicity, and the wave system are not specified.

In the usual unweighted scalar interpretation, Kirchhoff vertex conditions mean continuity at each vertex together with conservation of outward derivatives:
\[
 u_e(v,t)=u_{e'}(v,t)\quad\text{for all incident edge ends }e,e',
 \qquad
 \sum_{e\sim v}\partial_{\nu_e}u_e(v,t)=0.
\]
Here \(\partial_{\nu_e}\) denotes differentiation from the vertex into the incident edge, and the traces and sum must be well defined. The supplied record does not specify whether these conditions apply componentwise to a system, whether weights occur, or how boundary vertices are treated.

\textbf{Problem as supplied.} Can the nonlinear Schr\"odinger approximation be justified for the original nonlinear wave systems on periodic graphs with quadratic nonlinearities and Kirchhoff vertex conditions?

A justification would require a theorem comparing solutions of a specified original wave system with an approximation constructed from solutions of a specified NLS envelope equation. The record supplies neither equation, the admissible initial data, the carrier modes, the small parameter and scaling, the approximation formula, the error norm, nor the time interval and error bound. Consequently, it does not determine a fully quantified mathematical proposition. These missing choices must be recovered from the intended source before a precise approximation theorem can be stated; no particular choices are imposed here.

The supplied dated curation describes the problem as open and attributes the difficulty in the quadratic case to vertex regularity obstructing the required normal-form analysis. This is a reported source assessment, not an independently verified mathematical conclusion.
\subsection{Short English statement}
The problem asks whether NLS envelope dynamics rigorously approximate nonlinear wave systems with quadratic nonlinearities on periodic graphs satisfying Kirchhoff vertex conditions. The original system and the required approximation estimate are unspecified, so the supplied record supports a research question but not a fully quantified theorem statement. Its open status and explanation of the obstruction are supplied dated curation.
\subsection{Sources}
\begin{itemize}
\item Supplied local UnsolvedMath registry record, permanent ID \(30003521\), title \emph{NLS Approximation on Periodic Nonlinear Wave Graphs}, in \emph{Oberwolfach Reports Open Problems}; exact source problem code \texttt{catalogue entry}. The record gives proposed year 2017 and names no individual proposer.
\item Supplied report citation: \emph{Nonlinear Partial Differential Equations on Graphs (2018)}, Oberwolfach Reports, pp.~1805--1868, DOI \texttt{10.4171/owr/2017/29}. \url{https://doi.org/10.4171/owr/2017/29}. The background describes Oberwolfach Reports as documenting workshops and their open-problem discussions.
\item Supplied literature reference: Gilg--Pelinovsky--Schneider, NoDEA. \url{https://doi.org/10.1007/s00030-016-0417-7}.
\item Supplied literature reference: Gilg et al., Mathematische Nachrichten. \url{https://doi.org/10.1002/mana.202100009}.
\item Supplied literature reference: Oberwolfach Report 29/2017. \url{https://doi.org/10.4171/owr/2017/29}.
\item Literature assessment dated exactly \texttt{2026-08-21}: the supplied curation labels the problem open, states that the cubic NLS approximation on periodic quantum graphs is justified, and states that the quadratic-nonlinearity case remains explicitly open because vertex regularity obstructs the required normal-form analysis. This is supplied dated curation derived from the OWR record, not independently verified here; its time-sensitive status and cited claims require independent verification.
\end{itemize}
\end{document}
